\documentclass[11pt,a4paper]{article}

\usepackage[utf8]{inputenc}
\usepackage{geometry}
\geometry{a4paper, margin=1in}
\usepackage{graphicx}
\usepackage{booktabs}
\usepackage{hyperref}
\usepackage{amsmath}
\usepackage{multirow}
\usepackage{xcolor}
\usepackage{titlesec}

% Make section titles look like "1) Introduction"
\titleformat{\section}{\normalfont\Large\bfseries}{\thesection)}{1em}{}

\renewcommand{\contentsname}{Table of Contents}
\usepackage{tocloft}
\renewcommand{\cftsecaftersnum}{)}
\setlength{\cftsecnumwidth}{2em}

\usepackage{siunitx}
\sisetup{
    group-separator={,},
    group-minimum-digits=4,
    detect-weight=true,
    table-format=5.2
}

\title{\textbf{Modular Transformer Architecture for Language Modeling} \\ \large SAiDL Summer Induction Assignment}
\author{}
\date{May 2026}

\begin{document}

\maketitle

\tableofcontents
\newpage

\begin{abstract}
This report details the implementation, evolution, and evaluation of a modular Transformer language model trained on the WikiText-2 dataset. Our work progresses from a standard baseline architecture to advanced variants, incorporating alternative attention mechanisms, context-stable positional encodings, and local-global convolutional hybrids. Through systematic benchmarking, we identify architectures that optimize for inference speed, sequence length extrapolation, and computational efficiency.
\end{abstract}

\section{Introduction}
The Transformer architecture has revolutionized natural language processing, but its standard formulation suffers from quadratic computational complexity with respect to sequence length and struggles with out-of-distribution sequence length extrapolation. This project aims to iteratively build and evaluate a modular Transformer that addresses these bottlenecks. We focus on decoupling the core components---attention and positional encoding---and exploring state-of-the-art hybrid designs that integrate 1D convolutions to enhance local feature extraction and overall parameter efficiency.

\section{Methodology}

The project was executed in four sequential phases, progressively replacing standard components with more efficient or stable alternatives.

\subsection{Task 1: Baseline Architecture}
We established a standard autoregressive Transformer Language Model using absolute positional embeddings and standard $O(N^2)$ scaled dot-product attention. This model served as the foundational baseline to evaluate all subsequent architectural changes on the WikiText-2 dataset.

\subsection{Task 2: Attention Variants}
To address the quadratic bottleneck of standard attention, we implemented three highly efficient variants:
\begin{itemize}
    \item \textbf{Multi-Query Attention (MQA)}: Reduces memory bandwidth by sharing Key and Value heads across all Query heads.
    \item \textbf{Sliding Window Attention}: Restricts the attention receptive field to a localized window, enforcing $O(N)$ memory scaling for long documents.
    \item \textbf{Linear Attention}: Bypasses the $N \times N$ attention matrix entirely by utilizing a cumulative-sum formulation of the Key-Value interactions.
\end{itemize}

\subsection{Task 3: Positional Embedding and Context Extrapolation}
Absolute positional encodings fail mathematically when evaluating sequence lengths longer than the training context. To achieve true sequence extrapolation, we transitioned to relative positional strategies:
\begin{itemize}
    \item \textbf{Rotary Positional Embedding (RoPE)}: Encodes absolute positions with a rotation matrix and inherently introduces relative spatial distances into the dot-product attention.
    \item \textbf{Attention with Linear Biases (ALiBi)}: Removes positional embeddings entirely in favor of a static, non-learned distance penalty applied directly to the attention logits.
\end{itemize}

\paragraph{A Note on Positional Interpolation:}
Although the assignment optionally suggests implementing Positional Interpolation (scaling RoPE frequencies to fit longer context windows during evaluation), we did not implement this extension. Since standard RoPE and ALiBi successfully completed the sequence extrapolation tests (generalizing up to $L_{\text{test}}=2048$ tokens without numerical failure), interpolation was not necessary to meet the task's performance goals.


\subsection{Task 4: Convolution + Attention Hybrids}
To enhance the model's ability to capture short-range syntactic structures (n-grams) before engaging global attention, we implemented Causal 1D Convolutions using two designs:
\begin{itemize}
    \item \textbf{Design A (Pre-Attention)}: A Causal 1D Convolution acts as a local filter prior to the global attention mechanism.
    \item \textbf{Design B (Interleaved)}: Replaces alternating attention layers with Depthwise Separable Causal 1D Convolutions, drastically reducing parameter count and memory usage.
\end{itemize}

\subsection{Bonus Task: Attention-Free Transformer (AFT)}
To completely eliminate the $O(N^2)$ dot-product attention matrix, we integrated the Attention-Free Transformer (AFT) architecture. AFT utilizes a query-gated, element-wise aggregation with learned positional biases.

\begin{table}[h]
    \centering
    \caption{AFT Bonus Benchmarks}
    \begin{tabular}{lccc}
        \toprule
        \textbf{Variant} & \textbf{PPL} ($\downarrow$) & \textbf{Tokens/Sec} ($\uparrow$) & \textbf{VRAM (MB)} ($\downarrow$) \\
        \midrule
        AFT-Simple & 19298.07 & \textbf{29779.9} & \textbf{522.2} \\
        AFT-Local  & \textbf{17176.86} & 16803.3 & 590.1 \\
        AFT-Full   & 19405.19 & 16847.6 & 596.2 \\
        \bottomrule
    \end{tabular}
    \label{tab:aft_benchmarks}
\end{table}

The AFT variants dramatically shifted our performance landscape. \textbf{AFT-Local} surpassed all previous models to become the highest precision configuration, while \textbf{AFT-Simple} achieved the highest inference throughput of the entire project.

\section{Workflow and Implementation Requirements}

To ensure reproducibility across all project phases, we established a systematic execution workflow and identified the specific Python file requirements for each task. The central infrastructure consists of \texttt{model.py} (architecture), \texttt{config.py} (parameters), and \texttt{data.py} (WikiText-2 ingestion).

\begin{itemize}
    \item \textbf{Task 1 (Baseline)}: Requires \texttt{train.py}. \textit{Workflow}: Configure standard attention in \texttt{config.py} and execute \texttt{train.py}.
    \item \textbf{Task 2 (Efficiency)}: Requires \texttt{attention\_variants.py} and \texttt{evaluate\_all.py}. \textit{Workflow}: Toggle attention types in \texttt{config.py} and execute \texttt{evaluate\_all.py}.
    \item \textbf{Task 3 (Stability)}: Requires \texttt{positional\_logic.py} and \texttt{extrapolation\_test.py}. \textit{Workflow}: Set \texttt{pos\_type} to RoPE or ALiBi and execute \texttt{extrapolation\_test.py}.
    \item \textbf{Task 4 (Hybrids)}: Requires \texttt{conv\_logic.py} and \texttt{evaluate\_all.py}. \textit{Workflow}: Enable \texttt{use\_conv} in \texttt{config.py} and execute \texttt{evaluate\_all.py}.
    \item \textbf{Bonus Task (AFT)}: Requires \texttt{benchmark\_aft.py} and \texttt{attention\_variants.py}. \textit{Workflow}: Execute \texttt{benchmark\_aft.py}.
\end{itemize}

\section{Results and Comparison}

We conducted a global diagnostic across all 12 viable modular combinations, measuring Perplexity (PPL) on the validation set, Inference Speed (Tokens/Second), and Peak VRAM consumption (MB).

\subsection{Performance Topography}

\begin{table}[h]
    \centering
    \caption{Global Diagnostic Benchmarks (WikiText-2, Batch Size 4, Seq Len 128)}
    \begin{tabular*}{\textwidth}{@{\extracolsep{\fill}}l S[table-format=5.2] S[table-format=5.1] S[table-format=3.1] @{}}
        \toprule
        \textbf{Modular Variant} & {\shortstack{\textbf{PPL} \\ ($\downarrow$)}} & {\shortstack{\textbf{Tokens/Sec} \\ ($\uparrow$)}} & {\shortstack{\textbf{Peak VRAM} \\ (MB) ($\downarrow$)}} \\
        \midrule
        Standard + Absolute (Baseline) & 34940.81 & 12560.3 & 525.4 \\
        Standard + RoPE               & 29708.72 & 17744.2 & 525.4 \\
        Standard + ALiBi              & 29508.13 & 29206.4 & 525.4 \\
        \midrule
        MQA + Absolute                & 32572.44 & 30447.3 & 524.8 \\
        MQA + RoPE                    & 33442.34 & 26203.8 & 524.9 \\
        MQA + ALiBi                   & 30014.46 & 27564.7 & 524.8 \\
        \midrule
        Sliding Window + Absolute     & 31787.91 & 29841.4 & 525.4 \\
        Sliding Window + RoPE         & 29939.30 & 27420.6 & 525.4 \\
        Sliding Window + ALiBi        & 30421.29 & 28666.8 & 525.4 \\
        \midrule
        Linear + RoPE                 & 30210.36 & 16543.6 & 540.1 \\
        \midrule
        \textbf{HYBRID (Pre-Attn) MQA+ALiBi} & \textbf{27515.80} & 12075.8 & 526.5 \\
        \textbf{HYBRID (Interleaved) MQA+ALiBi} & 28808.97 & \textbf{24248.7} & \textbf{521.3} \\
        \bottomrule
    \end{tabular*}
    \label{tab:global_diagnostics}
\end{table}

\subsection{Sequence Length Extrapolation}
To evaluate length generalization, models trained on sequence length $L_{train}=512$ were evaluated at longer contexts ($L_{test}=1024, 2048$).

\begin{table}[h]
    \centering
    \caption{Extrapolation Stability Across Sequence Lengths}
    \begin{tabular}{lccc}
        \toprule
        \textbf{Positional Encoding} & $L_{test} = 512$ & $L_{test} = 1024$ & $L_{test} = 2048$ \\
        \midrule
        Absolute & 28,787 & 29,122 & \textcolor{red}{\textbf{FAILED (INF)}} \\
        RoPE & 30,510 & 30,864 & 31,489 \\
        ALiBi & 29,583 & 30,165 & 30,757 \\
        \bottomrule
    \end{tabular}
    \label{tab:extrapolation}
\end{table}

The Absolute Positional Encoding catastrophically failed when exposed to out-of-distribution sequence lengths, producing infinite perplexity. Conversely, both RoPE and ALiBi maintained robust mathematical stability, validating their efficacy in infinite-context or streaming scenarios.

\section{Conclusion}

This project successfully demonstrated that the standard Transformer architecture can be heavily optimized for speed, context length, and memory limits. 

\begin{itemize}
    \item \textbf{Efficiency \& Speed Winner}: \textbf{AFT-Simple} achieved the highest inference speed (nearly 30k Tok/Sec) while maintaining a minimal memory footprint (522 MB). For extreme memory-constrained edge deployments, the \textbf{Interleaved Convolutional Hybrid} remains highly competitive (521.3 MB).
    \item \textbf{Accuracy Winner}: \textbf{AFT-Local} achieved the lowest perplexity (17,176.86 PPL), significantly outperforming the Pre-Attention Convolutional Hybrid. By restricting the learned pairwise biases to a local sliding window, the model achieves superior grammar and pattern recognition without the noise of distant prior tokens.
\end{itemize}

By systematically decoupling the global sequence routing from localized processing, we engineered an adaptable and robust infrastructure capable of sustaining long-range dependencies without the crippling computational overhead of early sequence models.

\end{document}
